\documentclass[11pt]{article}

\usepackage[margin=1in]{geometry}
\usepackage{amsmath,amssymb,amsthm}
\usepackage{enumitem}
\usepackage{graphicx}
\usepackage{hyperref}

% Common PDE notation
\newcommand{\R}{\mathbb{R}}
\newcommand{\pd}[2]{\frac{\partial #1}{\partial #2}}
\newcommand{\pdd}[2]{\frac{\partial^2 #1}{\partial #2^2}}
\newcommand{\lap}{\Delta}
\newcommand{\grad}{\nabla}

% One numbered problem per \begin{problem}...\end{problem}
\theoremstyle{definition}
\newtheorem{problem}{Problem}

\title{APM5557: Partial Differential Equations\\Homework 2 Problems}
\author{Frank Palazzolo}
\date{}

\begin{document}

\maketitle

\begin{problem}
From problem 4.1-3 in the book

Prove that the solution of the Robin or third boundary value problem for the
Laplace equation:

$\lap u = 0$ in $\Omega$, with $\frac{\partial u}{\partial n} + \alpha u = \beta$ on $\partial \Omega$

is unique when $\alpha \ge 0$ is a constant.
\end{problem}

\begin{proof}[Solution]
Suppose $u_1$ and $u_2$ are two solutions to the Robin boundary value problem. Then their difference $v = u_1 - u_2$ satisfies the Laplace equation and the homogeneous Robin boundary condition:

$\lap v = 0$ in $\Omega$ and $\frac{\partial v}{\partial n} + \alpha v = 0$ on $\partial \Omega$.

Now we use Green's first identity for $v$ and $v$, which states that for a sufficiently smooth function $v$ on a domain $\Omega$ with boundary $\partial \Omega$:
\begin{equation} \label{eq:green-v}
\int_\Omega \left( v \lap v + |\grad v|^2 \right) dx
  = \int_{\partial\Omega} v \, \pd{v}{n} \, dS.
\end{equation}

Now, since $\lap v = 0$, the left-hand side simplifies to:
\begin{equation}
\int_\Omega |\grad v|^2 \, dx = \int_{\partial\Omega} v \, \pd{v}{n} \, dS.
\end{equation}

From the Robin boundary condition, we have $\pd{v}{n} = -\alpha v$ on $\partial \Omega$. Substituting this into the boundary integral gives:
\begin{equation}
\int_\Omega |\grad v|^2 \, dx = \int_{\partial\Omega} v (-\alpha v) \, dS = -\alpha \int_{\partial\Omega} v^2 \, dS.
\end{equation}

If $\alpha \ge 0$, then the right-hand side is non-positive, while the left-hand side is non-negative. 

This means that $\grad v = 0$ in $\Omega$ and $\alpha v = 0$ and $v = 0$ on $\partial \Omega$

Since $\grad v = 0$ in $\Omega$, $v$ is constant in $\Omega$, and since $v = 0$ on the boundary, it must be that $v = 0$ everywhere in $\Omega$.
\end{proof}

\begin{problem}
From problem 4.1-7 in the book

\end{problem}

\begin{proof}[Solution]
Solution goes here.
\end{proof}

\begin{problem}
From problem 4.1-9a in the book

\end{problem}

\begin{proof}[Solution]
Solution goes here.
\end{proof}

\end{document}
